\documentclass[11pt]{article}
\usepackage[a4paper,margin=25mm]{geometry}
\usepackage{amsmath,amssymb,amsthm,hyperref}
\hypersetup{colorlinks=true,linkcolor=blue,urlcolor=blue,citecolor=blue}
\newtheorem{theorem}{Theorem}
\newtheorem{lemma}[theorem]{Lemma}
\newcommand{\EE}{\mathbb E}
\newcommand{\PP}{\mathbb P}
\title{The two hyperbolic bounds for the Kilbas--Saigo function\\
\large MathDB 348998: proof throughout the conjectured parameter range}
\author{Principia Math In Progress}
\date{17 September 2026}
\begin{document}
\maketitle

\section*{Task 1: the full conjecture}
For $\alpha,m>0$ and $l>-1/\alpha$, write
\begin{equation}\label{definition}
 E_{\alpha,m,l}(z)=1+\sum_{n=1}^\infty
 \left(\prod_{k=1}^n
 \frac{\Gamma(1+\alpha((k-1)m+l))}
      {\Gamma(1+\alpha((k-1)m+l+1))}\right)z^n.
\end{equation}

\begin{theorem}\label{main}
Let $0<\alpha\leq1$, $m>0$, and $l>m-1/\alpha$. Define
\[
 c_- =\frac{\Gamma(1+\alpha(l-m))}
                 {\Gamma(1+\alpha(l-m+1))},
 \qquad
 c_+ =\frac{\Gamma(1+\alpha l)}
                 {\Gamma(1+\alpha(1+l))}.
\]
Then, for every $x\geq0$,
\begin{equation}\label{bounds}
 \boxed{\displaystyle
 \frac1{1+c_-x}\ \leq\ E_{\alpha,m,l}(-x)\ \leq\ \frac1{1+c_+x}.}
\end{equation}
Both inequalities are strict when $x>0$, and both are equalities at
$x=0$. The two constants are optimal within bounds of these forms.
\end{theorem}

This is Conjecture 4 in the published paper \cite{BS}, numbered
Conjecture 4.5 in its arXiv version. We prove both bounds with exactly
the printed hypotheses. The proof uses classical facts about
subordinators and their exponential functionals, credited to
\cite{BY}; the needed moment and comparison arguments are given below.

\subsection*{A positive gamma-ratio exponent}
Set
\begin{equation}\label{parameters}
 b=1+\alpha(l-m)>0,\qquad d=\alpha m>0,\qquad
 \phi(s)=\frac{\Gamma(b+ds+\alpha)}{\Gamma(b+ds)},\quad s\geq0.
\end{equation}
Thus $c_-=1/\phi(0)$ and $c_+=1/\phi(1)$.
For $0<\alpha<1$ the beta integral and integration by parts give
\begin{equation}\label{bernstein}
 \phi(s)=q+\frac{\alpha}{\Gamma(1-\alpha)}
 \int_0^\infty (1-e^{-dst})
       \frac{e^{-(b+\alpha)t}}{(1-e^{-t})^{1+\alpha}}\,dt,
 \qquad q=\phi(0)>0.
\end{equation}
Here is a direct derivation. For $v>0$,
\[
 \frac{\Gamma(v+\alpha)}{\Gamma(v)}
 =\frac{v}{\Gamma(1-\alpha)}
       \int_0^\infty e^{-vt}(e^t-1)^{-\alpha}\,dt
 =\frac{\alpha}{\Gamma(1-\alpha)}
       \int_0^\infty(1-e^{-vt})
       \frac{e^{-\alpha t}}{(1-e^{-t})^{1+\alpha}}\,dt.
\]
For the integration by parts, the boundary product
$(1-e^{-vt})(e^t-1)^{-\alpha}$ tends to zero at both ends:
it is $O(t^{1-\alpha})$ at zero and decays at infinity.
Subtract the last identity at $v=b$ from that at $v=b+ds$ to obtain
\eqref{bernstein}. All integrals so subtracted are finite.

The measure in \eqref{bernstein} has density of order
$t^{-1-\alpha}$ at zero and exponential decay at infinity.
After replacing the jump size $t$ by $dt$, it satisfies the
subordinator integrability condition $\int(1\wedge y)\,\nu(dy)<\infty$.
Consequently there exists a nondecreasing L\'evy process $S$, starting
at zero, such that
\begin{equation}\label{laplace}
 \EE e^{-sS_t}=e^{-t(\phi(s)-q)},\qquad s,t\geq0.
\end{equation}
This is the standard subordinator construction from its positive
jump measure; it also follows by truncating the small jumps of a
Poisson random measure and taking the increasing limit.
When $\alpha=1$, the same conclusion holds with
$\phi(s)=b+ds$, $q=b$, and the deterministic process $S_t=dt$.
Thus the endpoint $\alpha=1$ is covered separately and exactly.

\subsection*{Identifying the exponential functional}
Let $\zeta$ be independent of $S$, with exponential distribution of
rate $q$, and define
\[
 I=\int_0^\zeta e^{-S_t}\,dt.
\]
In particular $0<I\leq\zeta$ almost surely. For each integer $n\geq1$,
Tonelli's theorem and independent increments yield the classical
moment formula
\begin{equation}\label{moments}
 \EE I^n=\frac{n!}{\prod_{k=1}^n\phi(k)}.
\end{equation}
For completeness, expand the $n$th power as $n!$ times the integral
over $0<t_1<\cdots<t_n<\zeta$. Put $t_0=0$ and
$u_j=t_j-t_{j-1}$. The expectation of the integrand, including the
survival probability $e^{-qt_n}$, is
\[
 \exp\left(-\sum_{j=1}^n u_j\phi(n-j+1)\right).
\]
Integration over all positive $u_j$ proves \eqref{moments}. This is
also \cite[Theorem 2(i)]{BY}.

The function $\phi(s)$ increases to infinity. For $\alpha<1$ this
follows by monotone convergence in \eqref{bernstein}, since the
jump measure has infinite total mass; for $\alpha=1$ it is immediate.
Hence $\sum_{n\geq0}r^n/\prod_{k=1}^n\phi(k)$ converges for every
$r\geq0$ by the ratio test. Applying Tonelli to the exponential
series and then dominated convergence for complex arguments proves
\begin{equation}\label{representation}
 \EE e^{zI}=\sum_{n=0}^\infty\frac{z^n}{\prod_{k=1}^n\phi(k)}
            =E_{\alpha,m,l}(z),\qquad z\in\mathbb C.
\end{equation}
The last equality uses $b+dk=1+\alpha((k-1)m+l)$, so it identifies
every coefficient in \eqref{definition}, not merely an asymptotic
expansion. In particular
\[
 L(x):=E_{\alpha,m,l}(-x)=\EE e^{-xI},\qquad
 \mu:=\EE I=\frac1{\phi(1)}.
\]

\subsection*{The lower bound}
The pathwise inequality $I\leq\zeta$ gives, for $x\geq0$,
\begin{equation}\label{lower}
 L(x)\geq\EE e^{-x\zeta}=\frac{q}{q+x}
                         =\frac1{1+c_-x}.
\end{equation}

\subsection*{The upper bound from the remaining lifetime}
For $y>0$ let
\[
 A_t=\int_0^t e^{-S_s}\,ds,\qquad
 T_y=\inf\{t\geq0:A_t>y\}.
\]
The process $A$ is continuous and strictly increasing at finite
times. On $\{T_y<\infty\}$ one has $A_{T_y}=y$ and $T_y\geq y$.
Also $\{I>y\}=\{T_y<\zeta\}$ up to null events.
By the strong Markov property of $S$ and the memorylessness of the
independent killing time, the integral remaining after $T_y$, on
this event, has the form
\[
 I-y=e^{-S_{T_y}}I',
\]
where the future functional $I'$ has the same law as $I$ and is
independent of the past under the surviving conditional law.
Equivalently, without conditional notation,
\begin{align}
 \EE(I-y)_+
   &=\mu\,\EE\!\left[e^{-S_{T_y}-qT_y};\ T_y<\infty\right]
       \notag\\
   &\leq\mu\,\EE\!\left[e^{-qT_y};\ T_y<\infty\right]
    =\mu\,\PP(I>y).\label{residual}
\end{align}
The factor $e^{-qT_y}$ explicitly accounts for survival up to the
stopping time; it must not be omitted.

For $x>0$, integrate \eqref{residual} against $e^{-xy}\,dy$.
Tonelli applies to the nonnegative integrands. Direct integration
gives
\[
 \int_0^\infty e^{-xy}\EE(I-y)_+\,dy
       =\frac{\mu}{x}-\frac{1-L(x)}{x^2},\qquad
 \int_0^\infty e^{-xy}\PP(I>y)\,dy=\frac{1-L(x)}x.
\]
It follows that
\[
 \mu x-(1-L(x))\leq\mu x(1-L(x)),
 \qquad\text{and thus}\qquad
 L(x)\leq\frac1{1+\mu x}=\frac1{1+c_+x}.
\]
Together with \eqref{lower} and $L(0)=1$, this proves both conjectured
inequalities throughout the stated range.

\subsection*{Strictness, sharp constants, and an independent check}
For every fixed $t>0$, $S_t>0$ almost surely. For $\alpha=1$ this is
deterministic. For $\alpha<1$ it follows from
\[
 \PP(S_t=0)=\lim_{s\to\infty}\EE e^{-sS_t}
          =\lim_{s\to\infty}e^{-t(\phi(s)-q)}=0.
\]
Monotonicity of the paths and a countable set of positive rational
times give $S_t>0$ simultaneously for all $t>0$, almost surely.
Therefore $I<\zeta$ almost surely and \eqref{lower} is strict for
$x>0$. Since $I>0$ almost surely, there is an $\varepsilon>0$ such
that $\PP(I>\varepsilon)>0$. For every $0<y<\varepsilon$,
$T_y\geq y$ makes the inequality in \eqref{residual} strict.
Its integral is then strict, proving the asserted upper strictness.

For the constants, \eqref{definition} gives
$L(x)=1-c_+x+O(x^2)$ as $x\downarrow0$. A uniform upper bound
$L(x)\leq(1+cx)^{-1}$ therefore requires $c\leq c_+$.
The established asymptotic in \cite[Proposition 6]{BS} is
$L(x)\sim q/x$ as $x\to\infty$. A uniform lower bound
$(1+cx)^{-1}\leq L(x)$ consequently requires $c\geq1/q=c_-$.
This credits the source's already proved asymptotic; it is not a
new asymptotic claim here.

There is also a short independent check of the upper bound.
The classical factorization \cite[Theorem 2(ii)]{BY} supplies a
positive $R$, independent of $I$, with $IR$ standard exponential
and $\EE R=\phi(1)$. Conditional Jensen gives
\[
 L(x)\leq\EE\exp\left(-\frac{xIR}{\phi(1)}\right)
          =\frac1{1+x/\phi(1)}.
\]
The preceding remaining-lifetime proof establishes the bound
without relying on that factorization theorem.

\subsection*{Verification and source scope}
The verifier passes thirteen exact symbolic identities, including
parameter reindexing, the kernel integration by parts, the residual
Laplace calculation, and six beta-law moments at $\alpha=1$.
Ninety-six finite series diagnostics cover 24 parameter triples
at 160 and 220 digits. Twenty-four independent integrals check the
explicit $\alpha=1$ representation
$I=(1-e^{-d\zeta})/d$, and four quadratures check
\eqref{bernstein} at 100 digits. All checks pass. Numerical tests
are supplementary; the proof above applies to every allowed
parameter and every $x\geq0$.

The complete published source and its arXiv version retain the
conjecture. Dated bounded claim searches and a literature audit
accompany this note. Classical probability inputs are explicitly
credited. The authoring agent performed the audit; no independent
peer review is claimed.

\newpage
\section*{Tasks 2--5: proposed follow-up work}
These are four proposed tasks, not additional solved problems.
\begin{enumerate}
 \setcounter{enumi}{1}
 \item \textbf{Task 2.} Prove the conjectured sharp lower bound with
 power $1+1/m$ at the boundary $l=m-1/\alpha$, stated as
 Conjecture 3 in \cite{BS}, for every $m>0$.
 \item \textbf{Task 3.} Determine which changes of $\alpha,m,l$
 order the normalized variables $I/\EE I$ in convex order, and
 derive the corresponding parameter monotonicity of
 $E_{\alpha,m,l}(-x/\EE I)$.
 \item \textbf{Task 4.} Construct sharp rational refinements of
 \eqref{bounds} that match the coefficient of $x^2$ at zero and
 the leading $q/x$ term at infinity throughout the same range.
 \item \textbf{Task 5.} Characterize the parameters for which the
 survival function of $I$ is log-concave, strengthening the
 remaining-mean inequality \eqref{residual} to an increasing
 hazard rate.
\end{enumerate}

\begin{thebibliography}{9}
\bibitem{BS} L. Boudabsa and T. Simon,
\emph{Some Properties of the Kilbas--Saigo Function},
Mathematics \textbf{9}(3), 217 (2021), 24 pp.,
\href{https://doi.org/10.3390/math9030217}{doi:10.3390/math9030217}.
Conjecture 4 and Proposition 6, p.~16.
Preprint: \href{https://arxiv.org/abs/2012.05666}{arXiv:2012.05666},
Conjecture 4.5 and Proposition 4.3, p.~17.
\bibitem{BY} J. Bertoin and M. Yor,
\emph{Exponential functionals of L\'evy processes},
Probability Surveys \textbf{2} (2005), 191--212,
\href{https://doi.org/10.1214/154957805100000122}{doi:10.1214/154957805100000122}.
See \S3.1 and Theorem 2, pp.~193--196.
Complete journal reprint:
\href{https://arxiv.org/abs/math/0511265}{arXiv:math/0511265}.
\end{thebibliography}
\end{document}
